\section{Riesgos Éticos Identificados}

\subsection{Riesgo 1: Falta de Reproducibilidad Científica}

\subsubsection{Descripción del Riesgo}

Si los resultados de reducción de ruido no pueden ser reproducidos por investigadores independientes, la calidad científica se compromete. Esto es particularmente grave en astronomía, donde:

\begin{itemize}
    \item Las decisiones se basan en observaciones que pueden ser caras de replicar
    \item Los datos astronómicos son únicos en el tiempo (no se pueden retomar exactamente)
    \item La confianza pública en ciencia depende de rigor auditado
\end{itemize}

\subsubsection{Manifestación Potencial}

Un investigador reporta que su método reduce ruido con 95\% de exactitud, pero cuando otros intentan validarlo:
\begin{itemize}
    \item No pueden acceder al código fuente
    \item No saben exactamente qué datos fueron usados
    \item No tienen detalles sobre parámetros o configuración
    \item Los resultados reportados no pueden ser verificados
\end{itemize}

Resultado: La ciencia pierde credibilidad. Decisiones sobre observaciones futuras se basan en resultados no verificables.

\subsubsection{Mitigación Implementada}

Astro-Denoise implementa reproducibilidad mediante:

\begin{enumerate}
    \item \textbf{Código abierto}: Disponible en GitHub, con historial completo de versiones
    \item \textbf{Datos versionados}: Especifica exactamente qué versión de DC2 fue usada
    \item \textbf{Reportes estructurados}: Resultados en formato JSON/CSV auditables
    \item \textbf{Documentación técnica}: Guías paso-a-paso para replicar análisis
    \item \textbf{Cifras de validación cruzada}: Múltiples métricas independientes reportadas
\end{enumerate}

\subsection{Riesgo 2: Brecha de Acceso Institucional}

\subsubsection{Descripción del Riesgo}

Si solo investigadores en países aliados de LSST pueden validar y contribuir a ciencia astronómica moderna, se perpetúa una desigualdad sistémica que:

\begin{itemize}
    \item Limita oportunidades profesionales en el Sur Global
    \item Crea dependencia científica de instituciones del Norte
    \item Reduce diversidad en métodos y perspectivas de investigación
    \item Desaprovecha talento emergente
\end{itemize}

\subsubsection{Manifestación Potencial}

Un investigador en Bogotá o Lima identifica una mejora a un algoritmo de denoising, pero:
\begin{itemize}
    \item Su institución no es miembro de la alianza LSST
    \item Puede acceder a datos LSST, pero no a pipelines computacionales de referencia
    \item Publicar sobre métodos no estándar requiere más rigor de validación
    \item El circuito de revisión por pares favorece autores de instituciones prestigiosas
\end{itemize}

\subsubsection{Mitigación Implementada}

Astro-Denoise mitiga esto mediante:

\begin{enumerate}
    \item \textbf{Acceso sin restricción geográfica}: Código y datos sin barreras institucionales
    \item \textbf{Infraestructura mínima}: No requiere acceso a clusters especializados
    \item \textbf{Idioma accesible}: Documentación en español
    \item \textbf{Arquitectura modular}: Investigadores pueden extender con sus propios algoritmos
    \item \textbf{Estándares rigurosos}: Mismo nivel de validación para todos
\end{enumerate}

\subsection{Riesgo 3: Posible Falta de Gobernanza de Datos}

\subsubsection{Descripción del Riesgo}

Aunque Astro-Denoise usa datos simulados, la falta de políticas claras sobre quién puede usar los resultados, cómo se citan contribuciones, o qué pasa si datos se usan para fines no académicos, podría crear situaciones donde responsabilidades no son claras.

\subsubsection{Mitigación Implementada}

\begin{enumerate}
    \item \textbf{Licencias explícitas}: Código MIT, datos CC0 (dominio público)
    \item \textbf{Documentación de restricciones}: Uso permitido/no permitido claramente especificado
    \item \textbf{Compliance con normativas}: Respeto a Ley 1581
    \item \textbf{Atribución clara}: Cómo citar contribuciones científicas está documentado
\end{enumerate}

\subsection{Resumen de Riesgos y Controles}

\begin{table}[H]
\centering
\small
\begin{tabular}{|p{2.5cm}|p{3cm}|p{3.5cm}|}
\hline
\textbf{Riesgo} & \textbf{Severidad} & \textbf{Control Primario} \\
\hline
Irreproducibilidad & Alta & Código + datos abiertos, versionados \\
\hline
Acceso inequitativo & Media-Alta & Disponibilidad sin restricción geográfica \\
\hline
Gobernanza indefinida & Media & Licencias explícitas, documentación clara \\
\hline
\end{tabular}
\end{table}
